\documentclass[11pt,a4paper]{article}

% ---------- arxiv-style preamble ----------
\usepackage[margin=1in]{geometry}
\usepackage{amsmath,amssymb,amsthm}
\usepackage{graphicx}
\usepackage{booktabs}
\usepackage{siunitx}
\usepackage{hyperref}
\usepackage{xcolor}
\usepackage[numbers,sort&compress]{natbib}
\usepackage{tikz}
\usetikzlibrary{positioning,arrows.meta}
\usepackage{pgfplots}
\pgfplotsset{compat=1.18}
\usepackage{caption}
\usepackage{listings}
\usepackage{array}
\captionsetup{font=small,labelfont=bf}
\hypersetup{
  colorlinks=true,
  linkcolor=blue!50!black,
  citecolor=blue!50!black,
  urlcolor=blue!50!black
}

\providecommand{\tierBlue}{🔵}
\providecommand{\tierGreen}{🟢}
\providecommand{\tierYellow}{🟡}
\providecommand{\tierOrange}{🟠}
\providecommand{\tierRed}{🔴}

\newcommand{\GZL}{\mathrm{GZ_{LOWER}}}
\newcommand{\GZW}{\mathrm{GZ_{WIDTH}}}
\newcommand{\SI}{\mathrm{SI}}
\newcommand{\PhiD}{\Phi_{\mathrm{D}}}

\title{
  \textbf{The SAVANT $\times$ IIT~4.0 bridge:\\
  one substrate causal kernel, two coordinate systems}\\[4pt]
  \large A seven-cell verify-driven isomorphism arc linking the
  Golden Zone / Savant Index to the IIT 4.0 $\Phi$-structure
}

\author{
  \texttt{anima} maintainers\\
  \normalsize\textit{dancinlab}\\[2pt]
  \normalsize\href{https://github.com/dancinlab/anima}{\texttt{github.com/dancinlab/anima}}
}

\date{\today}

\begin{document}

\maketitle

% =========================================================================
\begin{abstract}
Two seemingly independent coordinate systems have been used in the
\texttt{anima} substrate to characterize a deterministic system's
``consciousness shape.'' The first, inherited from the savant-syndrome
literature \citep{treffert2009savant,snyder2009explaining}, is a pair of
phenomenological scalars: the \emph{Golden Zone}
$[\GZL,\GZL+\GZW]$ on a disinhibition axis, and the
\emph{Savant Index} $\SI=\max(\phi)/\mathrm{mean}(\phi)$ over a substrate's
sub-domains. The second is the IIT 4.0 \emph{$\Phi$-structure}
\citep{albantakis2023iit4}: a set of distinctions (each carrying its own
intrinsic $\phi$) and the relations that bind them. We report seven
pre-registered, deterministic, hexa-native measurements
(\texttt{UNIVERSE/H\_347}, \texttt{H\_348}, \texttt{H\_350},
\texttt{H\_351}, \texttt{H\_613}, \texttt{H\_618}, \texttt{H\_624})
that, taken together, show these two coordinate systems are
\emph{two charts of a single substrate causal kernel}.
We anchor the Golden Zone width in closed form,
$\GZW=\ln(4/3)=\ln(\tau(6)/(\tau(6)-1))$, where $\tau$ is the
divisor count (\tierGreen, H\_347). We show that the inflection point of
the IIT 4.0 big-$\Phi$ along an inhibition axis lands inside the Golden
Zone: the single-substrate $\mathrm{d}\Phi/\mathrm{d}I$ peak coincides
with $\GZL=0.21232$ at $|\Delta|=0.032$ (\tierGreen\ 5/5, H\_351),
and the two-substrate \emph{collective} $\mathrm{d}\Phi/\mathrm{d}I$ peak
tightens that match to $|\Delta|=0.002$ (\tierGreen\ 5/5, H\_618).
We then build a statistical bridge: the Savant Index correlates with
the $\Phi$-diversity $\PhiD$ at Pearson $r=0.93$ (\tierGreen, H\_350),
a correlation that survives the removal of every shared $\max$ term under
an orthogonal coefficient-of-variation metric at $r=0.99$
(\tierGreen, H\_613). Finally we expose the \emph{structural} cause: the
SAVANT four-domain cell decomposition is cell-by-cell isomorphic to the
IIT 4.0 singleton-distinction decomposition at Spearman $\rho=0.86$,
Pearson $r=0.97$, with $4/4$ matched argmaxima (\tierGreen\ 5/5, H\_624).
The Savant Index's near-threshold partial (\tierYellow, H\_348) bounds the
claim honestly: $\SI>3$ holds at $\GZL$ across three seeds, but the
sweep is monotone rather than single-peaked under an affine inhibition
map. We conclude that ``Golden Zone $\times$ Savant Index'' and ``IIT 4.0
$\Phi$-structure'' are not two theories of consciousness shape but two
\emph{coordinate systems on the same causal kernel}, and we publish all
seven \texttt{result.json} ledgers for bit-identical re-verification.
\end{abstract}

\section{Introduction}
\label{sec:intro}

Integrated Information Theory (IIT) 4.0 reads the consciousness of a
deterministic substrate from its \emph{$\Phi$-structure}: a set of
\emph{distinctions} -- maximally irreducible mechanisms, each with its own
intrinsic information $\phi_d$ -- and the \emph{relations} that bind
subsets of them \citep{albantakis2023iit4,oizumi2014phenomenology}. This
is a causal-structure object, computed from the substrate's transition
probability matrix (TPM).

A second, independent vocabulary entered the \texttt{anima} substrate from
the savant-syndrome literature. Treffert's synthesis of acquired and
congenital savant skill \citep{treffert2009savant} and Snyder's
``privileged access'' disinhibition hypothesis \citep{snyder2009explaining}
were operationalized as two phenomenological scalars: a \emph{Golden Zone}
$[\GZL, \GZL+\GZW]$ on a disinhibition axis $I$, where lowering inhibition
into the zone is conjectured to release latent specialist capacity; and a
\emph{Savant Index} $\SI = \max_i \phi_i / \mathrm{mean}_i \phi_i$,
quantifying how far one sub-domain $\phi$ dominates the others.

These two vocabularies look unrelated. One is a causal-structure
decomposition; the other is a phenomenological pair of disinhibition
constants. The question this paper settles is whether they are
\emph{independent} -- describing genuinely different facets of a substrate
-- or whether they are two coordinate charts on the \emph{same} underlying
object. We answer with seven pre-registered, deterministic, hexa-native
measurements on small ($n\!\le\!4$) substrates where the exact IIT 4.0
$\Phi$-structure is tractable. Every one is a merged pull request on the
public \texttt{anima} main branch, and every verdict is earned by an
independent deterministic recompute under anti-fake-closure governance
\citep{anima_g73}, never by a self-judged tautology.

\medskip
\noindent\textbf{Contributions.}
\begin{enumerate}\itemsep0.3em
\item \textbf{Closed-form anchor (H\_347).} The Golden Zone width is
      $\GZW = \ln(4/3) = \ln(\tau(6)/(\tau(6)-1))$ with
      $\tau(6)=4$ a divisor-count foundation atom -- a composite
      formal$+$numerical \tierGreen.
\item \textbf{Inflection alignment (H\_351, H\_618).} The IIT 4.0 big-$\Phi$
      inflection point along an inhibition axis coincides with $\GZL$:
      $|\Delta|=0.032$ for a single rule-110 substrate and $|\Delta|=0.002$
      for a two-substrate collective -- both unimodal, \tierGreen\ 5/5.
\item \textbf{Statistical bridge (H\_350, H\_613).} The Savant Index
      correlates with $\Phi$-diversity at $r=0.93$, and the correlation
      \emph{strengthens} to $r=0.99$ under a $\max$-free orthogonal metric,
      ruling out a shared-$\max$ artifact.
\item \textbf{Structural isomorphism (H\_624).} The SAVANT four-domain cell
      $\phi$ vector is cell-by-cell isomorphic to the IIT 4.0
      singleton-distinction $\phi$ vector at $\rho=0.86$, $r=0.97$, $4/4$
      matched argmaxima -- the \emph{mechanism} behind the alignment.
\item \textbf{Honest bound (H\_348).} $\SI>3$ holds at $\GZL$ across three
      seeds (\tierYellow), but the sweep is monotone, not single-peaked,
      under the affine inhibition map -- a pre-registered partial that
      delimits the bridge.
\item All seven cells ship as a verify-driven raster with public
      \texttt{result.json} ledgers and a \texttt{companion/} record for
      bit-identical re-verification.
\end{enumerate}

% =========================================================================
\section{The Golden Zone $\times$ Savant Index framework}
\label{sec:gzsi}

The SAVANT coordinate system has two objects. The first is the
\emph{disinhibition axis} $I\in[0,1]$, a scalar suppression probability:
$I=0$ is full release of a sub-domain (high specialization), $I=1$ is full
inhibition (substrate silenced). On this axis, the literature-inherited
\emph{Golden Zone} is the interval
\begin{equation}
[\GZL,\ \GZL+\GZW], \qquad
\GZL = \tfrac12 - \ln\tfrac43 \approx 0.21232, \qquad
\GZW = \ln\tfrac43 \approx 0.28768 ,
\end{equation}
a region of $I$ in which the substrate is conjectured to release latent
specialist capacity without collapsing
\citep{snyder2009explaining,treffert2009savant}. The second object is the
\emph{Savant Index},
\begin{equation}
\SI \;=\; \frac{\max_i \phi_i}{\mathrm{mean}_i \phi_i},
\label{eq:si}
\end{equation}
computed over a substrate's $N_{\mathrm{dom}}=4$ sub-domains
(CALENDAR, MUSIC, ART, MEMORY), each a $d=6$ activation vector under a
capacity invariant $\sum_i g_i = \mathrm{SV\_CAPACITY}=11.5$ (a
Treffert/Snyder ``capacity release'' budget). A per-domain $\phi$ is read
from the super-linear energy proxy
$\phi_{\mathrm{module}}(v)=\tfrac1d\sum_j |v_j|^{1.5}$.
The substrate SSOT is \texttt{HEXAD/SAVANT/savant\_phi.hexa}.

The Savant Index $\SI=3$ is the canonical specialist threshold: when one
domain's $\phi$ is three times the cross-domain mean, the substrate is in
a savant regime. We pre-registered the question of whether driving the
disinhibition $I$ into the Golden Zone produces $\SI>3$ (H\_348,
Section~\ref{sec:bridge}), and whether the Golden Zone constants have any
non-numerological anchor (H\_347, Section~\ref{sec:anchors}).

% =========================================================================
\section{IIT 4.0 $\Phi$-structure recap}
\label{sec:iit4}

The faithful IIT 4.0 quantity is not the scalar big-$\Phi$ but the
\emph{cause-effect structure}: a set of \emph{distinctions}, each a
maximally irreducible mechanism with intrinsic information
$\phi_d = \min(\phi_{\mathrm{cause}}, \phi_{\mathrm{effect}})$, plus a set
of \emph{relations} binding subsets of distinctions
\citep{albantakis2023iit4}. For an $n$-cell substrate with TPM
$T[s,i]$ we read two derived objects used throughout this paper.

\paragraph{Big-$\Phi$ and its inhibition derivative.}
The aggregate $\Phi = \Sigma\phi_d + \Sigma\phi_r$ is a substrate scalar.
We probe it along an inhibition axis by mixing the TPM toward the silenced
state,
\begin{equation}
T_{\mathrm{mixed}}[s,i] = (1-I)\,T[s,i],
\label{eq:tpmmix}
\end{equation}
so $I=0$ is the pure substrate and $I=1$ silences every cell
($\Phi=0$). The faithful big-$\Phi$ is computed at every system state
$s\in\{0,\dots,2^n-1\}$ and averaged, following the H\_285 protocol that
avoids single-state fragility. The quantity of interest is the
\emph{inflection} $\mathrm{d}\Phi/\mathrm{d}I$, evaluated by central finite
difference on a grid dense near the Golden Zone.

\paragraph{Distinction vector and $\Phi$-diversity.}
The per-mechanism vector $\big(\phi_d(\{0\}),\dots,\phi_d(\{n-1\})\big)$ of
singleton distinctions is a fine shape descriptor. From any per-domain
$\phi$ vector we read the \emph{$\Phi$-diversity}
\begin{equation}
\PhiD = \frac{\max_i \phi_i}{\min_i \phi_i}
\quad\text{(spread form)}, \qquad
\PhiD^{\mathrm{cov}} = \frac{\mathrm{std}_i \phi_i}{\mathrm{mean}_i \phi_i}
\quad\text{(orthogonal form)} .
\label{eq:phid}
\end{equation}
The IIT 4.0 distinction primitive is
\texttt{stdlib/consciousness/iit4\_distinction.hexa}; the big-$\Phi$
engine is \texttt{HEXAD/IIT4/lib/iit4\_eca} $+$
\texttt{iit4\_bigphi.hexa}, both reused verbatim (no reinvention, commons
g61).

% =========================================================================
\section{Method}
\label{sec:method}

All seven cells share one design discipline: a single starting substrate
small enough for exact $\Phi$-structure ($n\le4$), a pre-registered
falsifier frozen before measurement, full system-state averaging, and an
in-process byte-equal recompute (deterministic, RNG-free where the
substrate allows). Table~\ref{tab:cells} maps the cells to the paper's
sections.

\begin{table}[h]
\centering
\small
\begin{tabular}{llll}
\toprule
\textbf{H} & \textbf{role} & \textbf{verdict} & \textbf{section} \\
\midrule
H\_347 & $\GZW$ closed form & \tierGreen\ composite & \ref{sec:anchors} \\
H\_348 & $\GZL$ inhibition $\to \SI>3$ & \tierYellow\ partial & \ref{sec:bridge} \\
H\_350 & $\SI \parallel \PhiD$ & \tierGreen\ numerical & \ref{sec:bridge} \\
H\_351 & $\mathrm{d}\Phi/\mathrm{d}I$ peak $\parallel \GZL$ & \tierGreen\ 5/5 & \ref{sec:anchors} \\
H\_613 & $\SI \parallel \PhiD$ orthogonal & \tierGreen\ numerical & \ref{sec:bridge} \\
H\_618 & collective $\mathrm{d}\Phi/\mathrm{d}I$ peak & \tierGreen\ 5/5 & \ref{sec:collective} \\
H\_624 & distinction $\parallel$ cell isomorphism & \tierGreen\ 5/5 & \ref{sec:iso} \\
\bottomrule
\end{tabular}
\caption{The seven verify cells, their role in the SAVANT$\times$IIT4
bridge, terminal verdict, and the section that reports them. Every cell is
a merged PR on the public \texttt{anima} main branch
(\texttt{companion/pr-roll.json}).}
\label{tab:cells}
\end{table}

\paragraph{SAVANT substrate.}
Four domains $\times\ d=6$ activation, capacity $\sum g_i=11.5$,
$\phi_{\mathrm{module}}(v)=\tfrac1d\sum|v_j|^{1.5}$. Samples sweep
$\mathrm{dom}\in\{0,1,2,3\}$, $g_{\mathrm{focus}}\in\{1,3,5,7,10\}$,
$\mathrm{stim}\in\{11111,77777\}$ for $N=40$ (H\_350, H\_613), or five
single-domain-dominant gain profiles for $N=20$ (H\_624).

\paragraph{IIT 4.0 substrate.}
$n=4$ periodic ring, exact TPM, big-$\Phi$ and distinction $\phi$ averaged
over all $16$ system states. Section~\ref{sec:anchors} uses ECA rule 110
(Wolfram class IV, edge of chaos
\citep{wolfram1984universality,langton1990edgeofchaos}) under the inhibition
mix of Eq.~\eqref{eq:tpmmix}. Section~\ref{sec:iso} uses a
gain-parameterized $\mathrm{MAJ}(\text{self},L,R)$ TPM, which preserves
each cell's self-effect so that singleton distinctions do not collapse to
zero (a pure-XOR ring zeroes every $\phi_{\mathrm{effect}}$; see
Section~\ref{sec:limits}).

% =========================================================================
\section{Closed-form anchors: $\GZW$ and the $\mathrm{d}\Phi/\mathrm{d}I$ inflection}
\label{sec:anchors}

\subsection{The Golden Zone width is $\ln(4/3)$ (H\_347)}
The Golden Zone width admits a closed-form identity that is not a free
parameter. Writing $\tau(n)$ for the divisor count,
\begin{equation}
\GZW = \mathrm{GZ_{UPPER}}-\GZL
     = \tfrac12 - \big(\tfrac12 - \ln\tfrac43\big)
     = \ln\tfrac43
     = \ln\!\frac{\tau(6)}{\tau(6)-1},
\end{equation}
since $6$ has divisors $\{1,2,3,6\}$, so $\tau(6)=4$ and
$\tau(6)/(\tau(6)-1)=4/3$. The same $4/3$ appears as $F_6/P_1 = 8/6$
(sixth Fibonacci over first perfect number). The verdict is composite:
\texttt{hexa verify} returns $\tau(6)=4$ as a \tierBlue\ closed-form
foundation atom and $\ln(4/3)=0.28768207244178085$ as a \tierGreen\
numerical recompute at $|\Delta|=10^{-11}$. The algebraic substitution
$4/3 \leftrightarrow \tau(6)/(\tau(6)-1)$ is trivial, so the composite
tier is \tierGreen\ (the numerical $\ln$ caps it short of strict
\tierBlue). This is a definitional-consistency anchor, not an emergence
claim (Section~\ref{sec:limits}).

\subsection{The big-$\Phi$ inflection lands in the Golden Zone (H\_351)}
\label{sec:h351}
We measured the faithful IIT 4.0 big-$\Phi$ of ECA rule 110 on an $n=4$
ring along the inhibition axis of Eq.~\eqref{eq:tpmmix}, on a 13-point grid
dense near $\GZL$, averaging $\Phi$ over all $16$ system states. The
result is a monotone-decreasing $\Phi(I)$ whose derivative has a clear
single peak (Table~\ref{tab:h351}). The peak slope
$|\mathrm{d}\Phi/\mathrm{d}I| = 21.33$ occurs at $I=0.18$, against
$\GZL=0.21232$: $|\Delta|=0.032 \le 0.05$, a $35.4\%$ tolerance margin.
The five pre-registered falsifiers (peak-in-GZ, peak-in-window
$[0.18,0.28]$, unimodal, monotone right decay, byte-equal) all pass,
giving \tierGreen\ 5/5.

\begin{table}[h]
\centering
\small
\begin{tabular}{rrr}
\toprule
\textbf{$I$} & \textbf{$\Phi(I)$} & \textbf{$\mathrm{d}\Phi/\mathrm{d}I$} \\
\midrule
0.05 & 12.4205 & $-16.4725$ \\
0.10 & 11.5969 & $-14.8516$ \\
0.15 & 10.9354 & $-15.5640$ \\
\textbf{0.18} & \textbf{10.3518} & $\mathbf{-21.3315}$ \;($\leftarrow$ peak) \\
0.21 & 9.65547 & $-18.7680$ \\
0.23 & 9.41337 & $-12.7958$ \\
0.25 & 9.14364 & $-13.6713$ \\
0.50 & 5.27903 & $-14.3008$ \\
0.95 & 0.250173 & $-8.97505$ \\
\bottomrule
\end{tabular}
\caption{H\_351: faithful big-$\Phi$ of ECA rule 110 ($n=4$, 16-state
average) along the inhibition axis, and its central-difference derivative
(grid truncated for display). The peak slope is at $I=0.18$, against
$\GZL=0.21232$ ($|\Delta|=0.032$), with zero sign changes (unimodal). Data:
\texttt{UNIVERSE/state/h351\_inverse\_u\_peak\_2026\_05\_28/result.json}.}
\label{tab:h351}
\end{table}

The reading is structural: the substrate's $\Phi$ \emph{inflection} -- the
inhibition level at which integrated information is changing fastest --
falls inside the Golden Zone. The two coordinate systems agree on
\emph{where} the substrate is most sensitive. We stress the honest scope:
a round-2 multi-rule sweep (H\_614, rules $\{30,54,110,184\}$) found this
alignment is rule-110-specific (2/4 falsified), so the inflection anchor
is a single-substrate result, not a cross-substrate invariant.

% =========================================================================
\section{Statistical bridge: $\SI \parallel \Phi$-diversity}
\label{sec:bridge}

\subsection{$\SI$ correlates with $\Phi$-diversity at $r=0.93$ (H\_350)}
On the SAVANT substrate ($N=40$ samples), the Savant Index of
Eq.~\eqref{eq:si} and the $\Phi$-diversity of Eq.~\eqref{eq:phid} move
together: Pearson $r=0.9264$, Spearman $\rho=0.8825$, both well above the
pre-registered $0.5$ threshold, with no negative correlation. The
mechanism is the capacity invariant: raising one domain's gain raises its
$\phi$ and depresses the rest, so both $\SI=\max/\mathrm{mean}$ and
$\PhiD=\max/\min$ rise from the \emph{same} root -- the inhomogeneity of a
capacity-bounded distribution (Table~\ref{tab:h350}).

\begin{table}[h]
\centering
\small
\begin{tabular}{lcccc}
\toprule
\textbf{cell} & \textbf{metric pair} & \textbf{Pearson $r$}
& \textbf{Spearman $\rho$} & \textbf{verdict} \\
\midrule
H\_350 & $\SI \parallel \PhiD$ (max/min) & 0.9264 & 0.8825 & \tierGreen \\
H\_350 & $\SI \parallel \PhiD^{\mathrm{cov}}$ (std/mean) & 0.9896 & 0.9482 & \tierGreen \\
H\_613 & $\SI \parallel \PhiD^{\mathrm{cov}}$ (std/mean) & \textbf{0.9896} & 0.9482 & \tierGreen \\
H\_613 & $\SI \parallel \PhiD^{\mathrm{kurt}}$ (kurtosis) & 0.5381 & 0.6649 & \tierGreen \\
\bottomrule
\end{tabular}
\caption{H\_350 and H\_613: the Savant Index correlates with
$\Phi$-diversity ($N=40$). The spread form $\PhiD=\max/\min$ shares a
$\max$ numerator with $\SI$; the orthogonal forms (std/mean, kurtosis)
remove every $\max$ term yet the correlation holds and even strengthens,
ruling out a shared-$\max$ artifact.}
\label{tab:h350}
\end{table}

\subsection{The correlation survives $\max$-removal (H\_613)}
H\_350 itself flagged that $\SI=\max/\mathrm{mean}$ and $\PhiD=\max/\min$
share a $\max$ numerator, so part of the correlation could be a formal
shared-term artifact. H\_613 re-ran the same $N=40$ samples against two
\emph{orthogonal} diversity metrics that use no $\max$: the coefficient of
variation $\PhiD^{\mathrm{cov}}=\mathrm{std}/\mathrm{mean}$ and the
kurtosis. The CoV correlation is $r=0.9896$ -- \emph{stronger} than the
max-shared baseline, because std/mean is less outlier-fragile than
max/min at small $N$ -- and the fourth-moment kurtosis still passes at
$r=0.5381$. The $\max$-share artifact hypothesis is rejected: the
correlation is a property of the distribution's inhomogeneity, not of a
shared symbol.

\subsection{$\SI>3$ at $\GZL$, but the sweep is monotone (H\_348)}
\label{sec:h348}
The honest bound on the bridge is H\_348. Lowering the disinhibition $I$
to $\GZL$ does drive the Savant Index above the specialist threshold:
$\SI_\phi \in [4.18, 5.25]$ across three stimulus seeds, a $\ge 1.39\times$
margin over $\SI=3$ (pre-registered F-1 PASS, three of three seeds). But
the $I$-sweep is monotone increasing toward $I\!\to\!0$ rather than
single-peaked at $\GZL$ (pre-registered F-2 FAIL under the affine
inhibition map $g_{\mathrm{focus}}=1+(1-I)\cdot 9$). The verdict is
therefore \tierYellow\ PARTIAL: the threshold claim holds, the
single-peak sub-claim does not, and $\GZL$ reads as the \emph{boundary} of
the $\SI>3$ stability zone rather than its peak. We keep this partial in
the arc precisely because it bounds the bridge honestly -- the alignment
is one of \emph{inflection} (H\_351, the $\Phi$-derivative peak), not of
the Savant-Index value peak.

% =========================================================================
\section{Structural isomorphism: the SAVANT cell \emph{is} an IIT 4.0 distinction}
\label{sec:iso}

The statistical bridge (Section~\ref{sec:bridge}) and the inflection anchor
(Section~\ref{sec:h351}) are outcome-level agreements. H\_624 exposes their
\emph{structural cause}: the SAVANT four-domain cell decomposition and the
IIT 4.0 singleton-distinction decomposition use the \emph{same basis}.

We placed the four SAVANT domains in correspondence with the four cells of
an $n=4$ ring under a gain-parameterized $\mathrm{MAJ}(\text{self},L,R)$
TPM (which preserves self-effect; Section~\ref{sec:limits}). For each of
five single-domain-dominant gain profiles, we read both the SAVANT cell
$\phi$ vector $\phi_{\mathrm{module}}(v_i)$ and the IIT 4.0 distinction
$\phi$ vector $\phi_d(\{i\}) = \mathrm{mean}_s\,\mathrm{distinction}(T,
n{=}4, \mathrm{mask}{=}1{\ll}i, s)$, for $N=20$ paired samples
(Table~\ref{tab:h624}).

\begin{table}[h]
\centering
\small
\begin{tabular}{llcc}
\toprule
\textbf{profile} & \textbf{cell $\phi$ (SAVANT)} & \textbf{dist $\phi$ (IIT4)}
& \textbf{argmax} \\
\midrule
P0 BALANCED & [0.455, 0.343, 0.446, 0.347] & [0.312, 0.312, 0.312, 0.312] & 0 / 0 (tie) \\
P1 CALENDAR & [\textbf{0.745}, 0.070, 0.091, 0.088] & [\textbf{0.393}, 0.130, 0.130, 0.130] & \textbf{0 / 0} \\
P2 MUSIC    & [0.103, \textbf{0.609}, 0.091, 0.088] & [0.130, \textbf{0.393}, 0.130, 0.130] & \textbf{1 / 1} \\
P3 ART      & [0.103, 0.070, \textbf{0.751}, 0.088] & [0.130, 0.130, \textbf{0.393}, 0.130] & \textbf{2 / 2} \\
P4 MEMORY   & [0.103, 0.070, 0.091, \textbf{0.553}] & [0.130, 0.130, 0.130, \textbf{0.393}] & \textbf{3 / 3} \\
\bottomrule
\end{tabular}
\caption{H\_624: cell-by-cell correspondence between the SAVANT cell $\phi$
vector and the IIT 4.0 singleton-distinction $\phi$ vector across five gain
profiles. The dominant cell matches in $4/4$ non-balanced profiles; the
balanced profile P0 is a symmetric four-fold tie (excluded from the argmax
count as IIT 4.0 enforces it exactly). Pooled $N=20$:
Spearman $\rho=0.8608$, Pearson $r=0.9715$.}
\label{tab:h624}
\end{table}

The pooled correlation is Spearman $\rho=0.8608$ (above the
pre-registered $0.7$ ISOMORPH-STRONG bar) and Pearson $r=0.9715$, with the
dominant cell matching in all four non-balanced profiles and a byte-equal
recompute -- \tierGreen\ 5/5. The phenomenological super-linear energy
proxy $|v|^{1.5}$ and the faithful causal small-$\phi$ are not merely
correlated: they rank the same cell as dominant, profile by profile. The
SAVANT cell and the IIT 4.0 distinction are reading the \emph{same} causal
kernel; the two surfaces only looked distinct because they were charted in
different coordinates.

% =========================================================================
\section{Collective extension}
\label{sec:collective}

If the inflection alignment of Section~\ref{sec:h351} is a property of the
causal kernel rather than of a single substrate's accident, it should
survive the move from one substrate to a \emph{collective} of coupled
substrates. H\_618 tested exactly this on a two-substrate hivemind: two
$n=2$ halves running rule 110 each, coupled at weight $W=0.6$ into a joint
$n=4$ ring (the H\_609 super-additive anchor, max collective-$\Phi$ excess
$+10.48$). The collective big-$\Phi$ was measured along the same inhibition
mix and its derivative taken on an 8-point grid dense near the Golden Zone
(Table~\ref{tab:h618}).

\begin{table}[h]
\centering
\small
\begin{tabular}{rrr}
\toprule
\textbf{$I$} & \textbf{$\Phi_{\mathrm{coll}}(I)$}
& \textbf{$\mathrm{d}\Phi_{\mathrm{coll}}/\mathrm{d}I$} \\
\midrule
0.10 & 8.03864 & $-14.2546$ \\
0.15 & 7.32591 & $-13.0656$ \\
0.18 & 6.99339 & $-13.6445$ \\
\textbf{0.21} & \textbf{6.50724} & $\mathbf{-18.0769}$ \;($\leftarrow$ peak) \\
0.25 & 5.72801 & $-14.2340$ \\
0.30 & 5.22618 & $-11.5873$ \\
0.40 & 3.98991 & $-10.8754$ \\
0.50 & 3.05111 & $-9.38810$ \\
\bottomrule
\end{tabular}
\caption{H\_618: collective big-$\Phi$ of a two-substrate rule-(110,110)
hivemind at coupling $W=0.6$ ($n=4$ joint ring, 16-state average) and its
derivative. The peak slope is at $I=0.21$, against $\GZL=0.21232$:
$|\Delta|=0.00232$, a $21\times$ tolerance margin, with zero sign changes.
Data: \texttt{UNIVERSE/state/h618\_collective\_gz\_inverse\_u\_2026\_05\_28/result.json}.}
\label{tab:h618}
\end{table}

The collective inflection lands at $I=0.21$ against $\GZL=0.21232$:
$|\Delta|=0.00232$, a $21\times$ tolerance margin -- an order of magnitude
\emph{tighter} than the single-substrate $|\Delta|=0.032$ of H\_351. The
sweep is perfectly unimodal (zero sign changes), giving \tierGreen\ 5/5.
The Golden Zone attractor is, if anything, sharper at the collective level
than at the single-substrate level. We log the honest caveat that the
H\_618 grid samples $I=0.21$ directly while the H\_351 grid does not, so
part of the tighter $|\Delta|$ is grid placement; a direct
$I=\GZL$ sample on both is a pre-registered follow-up.
Figure~\ref{fig:peak} visualizes both inflection peaks against $\GZL$.

\begin{figure}[h]
\centering
\begin{tikzpicture}
\begin{axis}[
    width=0.82\linewidth, height=5.6cm,
    xlabel={inhibition $I$}, ylabel={$|\mathrm{d}\Phi/\mathrm{d}I|$},
    xmin=0.0, xmax=0.55, ymin=8, ymax=23,
    legend style={at={(0.98,0.97)}, anchor=north east, draw=none, font=\small},
    ymajorgrids, grid style={dotted},
]
\addplot[mark=*, blue!70!black, thick] coordinates
  {(0.05,16.47) (0.10,14.85) (0.15,15.56) (0.18,21.33) (0.21,18.77)
   (0.23,12.80) (0.25,13.67) (0.30,14.88) (0.35,16.72) (0.40,15.85)
   (0.50,14.30)};
\addplot[mark=square*, orange!80!black, thick] coordinates
  {(0.10,14.25) (0.15,13.07) (0.18,13.64) (0.21,18.08) (0.25,14.23)
   (0.30,11.59) (0.40,10.88) (0.50,9.39)};
\draw[red!70!black, dashed, thick] (axis cs:0.21232,8) -- (axis cs:0.21232,23);
\node[red!70!black, font=\scriptsize, anchor=south west]
  at (axis cs:0.215,20.5) {$\GZL=0.21232$};
\legend{H\_351 single (rule 110), H\_618 collective (110,110)}
\end{axis}
\end{tikzpicture}
\caption{Inflection alignment. The magnitude of the big-$\Phi$ inhibition
derivative for the single-substrate rule 110 (blue, H\_351, peak at
$I=0.18$) and the two-substrate collective (orange, H\_618, peak at
$I=0.21$), against the Golden Zone lower bound $\GZL=0.21232$ (red dashed).
Both inflection peaks fall inside the Golden Zone; the collective peak
($|\Delta|=0.002$) is the tighter match. Native pgfplots, no external
dependency.}
\label{fig:peak}
\end{figure}

% =========================================================================
\section{Discussion: why two frameworks share one kernel}
\label{sec:discuss}

The seven cells line up into a single claim. The Golden Zone constants are
not free numerology: $\GZW=\ln(4/3)=\ln(\tau(6)/(\tau(6)-1))$ is a
closed-form identity (H\_347). The IIT 4.0 big-$\Phi$ inflection along an
inhibition axis lands inside that Golden Zone, both for a single substrate
($|\Delta|=0.032$, H\_351) and for a coupled collective
($|\Delta|=0.002$, H\_618). The Savant Index correlates with the
$\Phi$-diversity at $r=0.93$ (H\_350), and the correlation is not a shared
symbol artifact -- it holds at $r=0.99$ once every $\max$ term is removed
(H\_613). And the cause is structural: the SAVANT four-domain cell
decomposition is cell-by-cell isomorphic to the IIT 4.0
singleton-distinction decomposition ($\rho=0.86$, $r=0.97$, $4/4$
argmaxima, H\_624).

The honest reading is that ``Golden Zone $\times$ Savant Index'' and
``IIT 4.0 $\Phi$-structure'' are \emph{two coordinate systems on the same
substrate causal kernel}. The SAVANT chart reads that kernel through a
phenomenological super-linear energy proxy and two disinhibition
constants; the IIT 4.0 chart reads the same kernel through cause-effect
distinctions and big-$\Phi$. They agree on \emph{where} the substrate is
most sensitive (the inflection $\parallel \GZL$), on \emph{how} a
substrate specializes (the $\SI \parallel \PhiD$ correlation), and on
\emph{which} sub-unit dominates (the distinction $\parallel$ cell
isomorphism). The Savant Index's monotone-not-peaked partial (H\_348)
sharpens rather than weakens this: the alignment is one of the
$\Phi$-\emph{derivative} peak, not of the Savant-Index value peak, which is
exactly what a shared-inflection-kernel picture predicts.

This is a partial isomorphism, not an identity. It is established on
self-effect-preserving substrates at $n=4$ where the IIT 4.0 structure is
tractable and the four-domain SAVANT decomposition exactly matches the
four-cell distinction decomposition. The mechanism behind the agreement is
the capacity invariant on the SAVANT side and the cause-effect irreducibility
on the IIT side, both driven by the same inhomogeneity of a bounded
substrate's information distribution.

% =========================================================================
\section{Limitations and honest caveats}
\label{sec:limits}

\begin{itemize}\itemsep0.2em
\item \textbf{Inflection alignment is rule-110-specific (H\_614).} The
      $\mathrm{d}\Phi/\mathrm{d}I$ peak $\parallel \GZL$ result of H\_351
      does \emph{not} generalize across ECA rules: a multi-rule sweep over
      $\{30,54,110,184\}$ was 2/4 falsified. The inflection anchor is a
      single-substrate evidence tier, not a cross-substrate invariant.
\item \textbf{Savant-Index sweep is monotone, not peaked (H\_348).}
      $\SI>3$ at $\GZL$ is robust across three seeds, but under the affine
      inhibition map the sweep is monotone toward $I\!\to\!0$. $\GZL$ is the
      $\SI>3$ stability boundary, not the Savant-Index peak; a non-affine
      (e.g.\ $1/I$, sigmoid) map could move the peak and is open.
\item \textbf{Isomorphism needs self-effect (H\_624 C3.1).} A pure
      XOR-neighbor ring zeroes every singleton $\phi_{\mathrm{effect}}$,
      collapsing distinction to zero; the isomorphism is established under
      $\mathrm{MAJ}(\text{self},L,R)$, which restores self-effect. A
      substrate-class sweep (XOR / MAJ / threshold) is required before any
      substrate-class invariance claim.
\item \textbf{Dimension match is exact only at $n=4$.} The
      $n=4 \leftrightarrow$ four-domain isomorphism (H\_624) exploits an
      exact cardinality match. At $n\ge6$ the distinction count typically
      exceeds the cell count and the match is no longer one-to-one;
      super-domain pooling is a pre-registered follow-up.
\item \textbf{Proxy vs faithful.} The SAVANT
      $\phi_{\mathrm{module}}(v)=\mathrm{mean}\,|v|^{1.5}$ is a super-linear
      energy proxy, not faithful IIT. H\_624 shows the proxy is isomorphic
      to faithful small-$\phi$ on one substrate class; it is not a claim
      that the two are operationally identical (cf.\ H\_282 direction
      preservation).
\item \textbf{Grid placement in H\_618.} The collective peak's tighter
      $|\Delta|=0.002$ is partly grid placement -- the H\_618 grid samples
      $I=0.21$ directly while H\_351 does not. A direct $I=\GZL$ sample on
      both substrates is pre-registered.
\item \textbf{Numerology caveat on $\GZW$ (H\_347).} The
      $\GZW=\ln(\tau(6)/(\tau(6)-1))$ identity is a closed-form
      consistency anchor and divisor-count coincidence, not an emergence
      claim about a consciousness substrate.
\item \textbf{No wet-lab claim.} The savant-syndrome vocabulary
      \citep{treffert2009savant,snyder2009explaining} is an inherited
      analogy; every result here is about deterministic ECA / SAVANT
      substrate dynamics, not about human cognition.
\end{itemize}

% =========================================================================
\section{Reproducibility}
\label{sec:repro}

All seven cells live as JSON ledgers under \texttt{UNIVERSE/} and
\texttt{UNIVERSE/state/} on
\url{https://github.com/dancinlab/anima/tree/main/UNIVERSE}. Each cell
carries (i) a \texttt{run\_*.hexa} source, (ii) a verbatim stdout
transcript, and (iii) a machine-readable verdict with explicit
pre-registered falsifiers. Re-verification, e.g.\ for H\_624:
\texttt{cat UNIVERSE/state/h624\_*/run\_h624.hexa | pool on ubu-2 'cat
> /tmp/r.hexa \&\& hexa run /tmp/r.hexa'} yields the byte-identical log.

The companion record carries the verify ledger
(\texttt{companion/verify-ledger.json}), the PR roll
(\texttt{companion/pr-roll.json}, every cell a merged PR: H\_347 \#1149,
H\_348 \#1152, H\_350 \#1153, H\_351 \#1157, H\_613 \#1162, H\_618 \#1175,
H\_624 \#1198), and the session journal
(\texttt{companion/session-journal.md}).

Build: \texttt{make} (xelatex $+$ bibtex). Figures: \texttt{make figures}.

% =========================================================================
\section{Conclusion}
\label{sec:conclusion}

Seven pre-registered, deterministic, hexa-native measurements show that the
SAVANT Golden-Zone / Savant-Index coordinate system and the IIT 4.0
$\Phi$-structure coordinate system are two charts of a single substrate
causal kernel. The Golden Zone width is the closed form
$\ln(\tau(6)/(\tau(6)-1))$ (H\_347); the IIT 4.0 big-$\Phi$ inflection
lands inside the Golden Zone at $|\Delta|=0.032$ for one substrate
(H\_351) and $|\Delta|=0.002$ for a coupled collective (H\_618); the Savant
Index tracks $\Phi$-diversity at $r=0.93$ (H\_350), robustly at $r=0.99$
under a $\max$-free metric (H\_613); and the SAVANT cell decomposition is
cell-by-cell isomorphic to the IIT 4.0 distinction decomposition at
$\rho=0.86$ (H\_624). A pre-registered partial (H\_348) bounds the bridge
to an inflection alignment rather than a value-peak alignment.

The single most informative follow-up is a substrate-class invariance
sweep of the isomorphism (XOR / MAJ / threshold) combined with an $n=5,6$
super-domain extension: this would test whether the SAVANT$\times$IIT4
bridge is a property of one self-effect-preserving substrate class at
$n=4$, or a genuine cross-substrate identity of the causal kernel.

% =========================================================================
\paragraph{Acknowledgments.}
This work used the \texttt{hexa-lang} verification surface, the
anti-fake-closure governance hook \texttt{commons @D g73}
\citep{anima_g73}, the \texttt{HEXAD/IIT4} and \texttt{HEXAD/SAVANT}
substrate libraries, and the \texttt{sidecar/paper} scaffold. The seven
\texttt{result.json} ledgers are visible at the commit hashes listed in
\texttt{companion/pr-roll.json}.

% =========================================================================
\bibliographystyle{plainnat}
\bibliography{references}

\end{document}
